\documentclass[letterpaper,twoside,11pt]{book}

\usepackage{ustj} % Custom package for styling
% \usepackage[utf8]{inputenc} % Character encoding
% \usepackage{fancyhdr} % Custom headers/footers (already loaded in ustj, but explicit for clarity)
% \usepackage{bibentry} % For bibliography
\addbibresource{main.bib}

\usepackage{fontspec}
\setsansfont{Helvetica}
\newfontfamily\gothicfont{Textura Libera Tenuis}

\usepackage{tikz}
\usepackage{pgfplots}
\pgfplotsset{compat=1.10}
\usetikzlibrary{shapes.geometric,arrows,fit,matrix,positioning}
\tikzset
{
    treenode/.style = {circle, draw=black, align=center, minimum size=1cm},
    subtree/.style  = {circle, draw=darkgray, dashed, text=darkgray, align=center, minimum size=1cm}
}

% Define author commands
\newcommand{\authorname}{N.\ E.\ Davis}
\newcommand{\authorpatp}{\patp{lagrev-nocfep}}

% Define custom abstract environment for chapters
\newenvironment{abstract}
  {\par\noindent\textbf{Abstract}\par\nobreak\smallskip}
  {\par\bigskip}

% Define chaptertitle command for use in headers
\newcommand{\chaptertitle}{}

%% MANUSCRIPT
\title{\fontsize{36pt}{36pt}\gothicfont NOCKONOMICON}
\author{\fontsize{24pt}{24pt}\gothicfont\authorname}
\date{}

\begin{document}

\frontmatter %%%%%%%%%%%%%%%%%%%%%%%%%%%%%%%%%%%%%%%%%%%%%%%%%%%%%%%%%%%%%%%%%%%
\fancyhf{}
\fancyhead[LE]{}
\fancyhead[RO]{}
\fancyfoot[LE,RO]{}

% Title page
\maketitle

% Dedication and copyright page
\cleardoublepage
\thispagestyle{empty} % No page number on this page

% Table of contents
\tableofcontents

\mainmatter %%%%%%%%%%%%%%%%%%%%%%%%%%%%%%%%%%%%%%%%%%%%%%%%%%%%%%%%%%%%%%%%%%%%

\renewcommand{\chaptertitle}{The Roots of Computing in the Ancient World}

%  Make first page footer:
\fancypagestyle{firststyle}{%
\fancyhf{}% Clear header/footer
\fancyhead{}
\fancyfoot[L]{{\footnotesize
              %% We toggle between these:
              % Manuscript submitted for review.\\
              \authorname{}~\authorpatp{}.  ``\chaptertitle{}''.  {\it The Nock Book} (2025):  1–$x$.
              }}
}
%  Arrange subsequent pages:
\fancyhf{}
\fancyhead[LE]{{\urbitfont The Nock Book}}
\fancyhead[RO]{\chaptertitle{}}
\fancyfoot[LE,RO]{\thepage}

\chapter{\chaptertitle}

\thispagestyle{firststyle}

\begin{abstract}
Cuneiform and early computational concepts are explored in this chapter.
\end{abstract}

\begin{quote}
Machines are pure products of human intellect, constructed for man’s own purposes. They are nowhere to be found in the world of nature (as products of nature); yet the workings of the laws of nature find their purest expression in machines, purer than in any of the products of nature itself. The laws of nature work directly in machines, with an immediacy not to be found in the products of nature.

(Keiji Nishitani, {\it Religion and Nothingness}, 1982, pp.~129--130)
\end{quote}

There are many worthy raconteurs of the history of computing, such as \citet{Hofstadter1980} and \citet{Penrose2004}.  What distinguishes my approach to the history of computing from other tellings?  Computing has deep connexions to alchemy and true language, which I hope to demonstrate without assuming any Hermetic background on the reader's part.  Much like technology, computing is not merely a tool but a stance towards the world.  However, it does not entail technologism's supercession; computing is always both a lens and a kaleidoscope, through which many layers of reality can be seen simultaneously.  It is a way of seeing the world as a system of pure statements and acting to make those statements true.  It is both ontology and epistemology.

First, therefore, I must establish that computation is the apex of a long human search for a language of truth.  The origin of symbol itself in the σύμβολον, a token broken in half to infallibly mark identity, is one font.  Later ``symbol'' was the profession of faith, both an attempt to formulate truth and a sign of shared claims.  The primary source, however, is language itself, Hofstadter's ``strange loop'' of self-reference.  The first possible Nock formula is a crash, the void.  The second is self-reference, awakening.  Peirce would be proud of the progression.  We have tricked ourselves into thinking that computing is a late development, a fruitful elaboration of the end of mathematics.  In fact, it is the philosopher's stone, and it underlies every process in the world.  It has turned lead into gold.

The historical account which follows aims throughout at the goal of understanding what Nock is and why it matters as a protocol.  The reader will understand what mankind has been looking for, when we have caught various pieces of the puzzle, how these things came together in the mid-20th century, and how they have been refined into a pure machine.

Outline:
- Cuneiform, tabulation, calculation, and algorithm
- Solid-state indices:  abaci, tokens, etc.  Mechanism as embodied equation.
- Pure language and true statements (Llull, Leibniz, Wilkins)
- Generalization (Jacquard, Babbage)
- The universal machine (Turing, G\"{o}del, Zuse, von Neumann)
- Pure statements as Pages from the Book (Curry, Church)

\section{Tabulation}

\section{Solid-state indices}

The problem is that any calculation lives in someone's head or someone's claim.  This is a summary, but not the proof or solution.

\renewcommand{\chaptertitle}{Combinators for Fun and Profit}

%  Make first page footer:
\fancypagestyle{firststyle}{%
\fancyhf{}% Clear header/footer
\fancyhead{}
\fancyfoot[L]{{\footnotesize
              %% We toggle between these:
              % Manuscript submitted for review.\\
              \authorname{}~\authorpatp{}.  ``\chaptertitle{}''.  {\it The Nock Book} (2025):  10–$x$.
              }}
}
%  Arrange subsequent pages:
\fancyhf{}
\fancyhead[LE]{{\urbitfont The Nock Book}}
\fancyhead[RO]{\chaptertitle{}}
\fancyfoot[LE,RO]{\thepage}

\chapter{\chaptertitle}

\thispagestyle{firststyle}

\begin{abstract}

\end{abstract}

The first section of this chapter will introduce the basic \textct{SKI} and \textct{BCKW} combinator calculi and examine how they can represent each other.  The second section will introduce Nock's opcodes as combinators, and show how they can be used to express various combinator idioms.  The third section will show how to build practical affordances for common programming tasks, such as conditionals, recursion, subprocedure invocation, and data structures.  The final section will discuss the relationship between combinators and virtualization, and how Nock can be used to build a virtual machine even in a crash-only context.

\section{Combinators in Computing}

The pioneers of computing elaborated three complementary visions:

\begin{enumerate}
  \item  The Turing machine, a simple abstract machine that can compute anything computable.
  \item  The lambda calculus, a pure functional language based on substitution and abstraction.
  \item  Combinatory logic, a variable-free functional language based on combinators.
\end{enumerate}

Logicians early demonstrated the equivalence of these approaches, but particular formulations are still preferred for particular tasks.  The Turing machine is the most intuitive model of computation, and is often used to prove computability results.  The lambda calculus is the basis of many functional programming languages, such as Lisp, Scheme, and Haskell.  Combinatory logic is less widely used in practice, but has a long history in mathematical logic and type theory.  It is also the basis of Nock.

Before we get too far out over our skis, let's talk about what a combinator is.  At its simplest, a combinator is a function that takes one or more arguments and returns a result, without referring to any variables outside its own arguments.  In practice, we'll use different combinators to build up more complex functions, and we'll see informally how they can be combined to express any computable function.  We can even build combinators out of other combinators.  The objective of combinatory logic is to express all computation in terms of combinators alone, without the need for variables or named functions.

This first section will introduce some elemental and compound combinators and make the case that they circumscribe ``computing'' as a concept.

\subsection{The \textct{SKI} calculus}

Computable functions can be expressed using minimal sets of combinators.  Many logicians today prefer to work with Moses Schönfinkel and Haskell Curry's \textct{SKI} combinator calculus, which has three combinators:


The identity combinator \textct{I} is conventionally written \textct{Ix = x}.  This simply means that it returns its single argument unchanged.  There is little to say of this combinator at this point, but logicians have found that it frequently simplifies expressions.

\begin{lstlisting}[style=listingcomb]
I x
x
\end{lstlisting}

The constant combinator \textct{K} is conventionally written \textct{Kxy = x}.  This means that it takes two arguments and returns the first, ignoring the second.  It is useful for building functions that always return a fixed value.

\subsubsection{Questions}

\begin{enumerate}
  \item  What is the result of \textct{I x}?
  \item  What is the result of \textct{I I x}?
\end{enumerate}

\begin{lstlisting}[style=listingcomb]
K x y
x
\end{lstlisting}

The substitution combinator \textct{S} is conventionally written \textct{Sxyz = xz(yz)}.  This means that it takes three arguments: a function \textct{x}, and two values \textct{y} and \textct{z}.  It applies the function \textct{x} to the value \textct{z}, and also applies the function \textct{y} to the value \textct{z}, then applies the result of \textct{x z} to the result of \textct{y z}.  This combinator is more complex than the others, but it is very powerful, as it allows us to express function application and composition.

\begin{lstlisting}[style=listingcomb]
S x y z
x z (y z)
\end{lstlisting}

The \textct{SKI} system is Turing complete \citep{blah}, meaning that any computable function can be expressed using only these three combinators.  It is also equivalent to the lambda calculus, meaning that any lambda expression can be translated into an equi\-valent \textct{SKI} expression and vice versa.  (cite church-turing thesis stuff)

Combinators can be combined to form more complex expressions.  For example, the combinator \textct{S K K} is equivalent to the identity combinator \textct{I}:

\begin{lstlisting}[style=listingcomb]
S K K x
K x (K x)
x
\end{lstlisting}

% https://en.wikipedia.org/wiki/Magnum_opus_(alchemy)

% https://en.wikipedia.org/wiki/B,_C,_K,_W_system
% https://dkeenan.com/Lambda/index.htm

% https://doisinkidney.com/posts/2020-10-17-ski.html
% https://www.youtube.com/watch?v=GawiQQCn3bk

\subsection{The \textct{BCKW} calculus}

The \textct{BCKW} combinator calculus is an alternative to \textct{SKI}, with a different set of primitive combinators.  It has four combinators:

The composition combinator \textct{B} is conventionally written \textct{Bxyz = x(yz)}.  This means that it takes three arguments: a function \textct{x}, and two values \textct{y} and \textct{z}.  It applies the function \textct{y} to the value \textct{z}, then applies the function \textct{x} to the result.

\begin{lstlisting}[style=listingcomb]
B x y z
x (y z)
\end{lstlisting}

% The constant combinator \textct{C} is conventionally written \textct{Cxyz = xzy}.  This means that it takes three arguments: a function \textct{x}, and two values \textct{y} and \textct{z}.  It applies the function \textct{x} to the values \textct{z} and \textct{y}, effectively reversing the order of the last two arguments.

% The weakening combinator \textct{W} is conventionally written \textct{Wxy = xyy}.  This means that it takes two arguments: a function \textct{x} and a value \textct{y}.  It applies the function \textct{x} to the value \textct{y} twice.

% The duplication combinator \textct{U} is conventionally written \textct{Uxyz = xzyz}.  This means that it takes three arguments: a function \textct{x}, and two values \textct{y} and \textct{z}.  It applies the function \textct{x} to the values \textct{z}, \textct{y}, and \textct{z} again.  This combinator is less common, but can be useful in certain contexts.

The \textsc{SKI} and \textsc{BCKW} systems are equivalent, meaning that any expression in one system can be translated into an equivalent expression in the other system.  For example, the identity combinator \textct{I} can be expressed in \textsc{BCKW} as:

\begin{lstlisting}[style=listingcomb]
I = W K

\end{lstlisting}


\textct{U}

\textct{Y}

\begin{lstlisting}[style=listingcomb]
Y = S (K (S I I)) (S (S (K S) K) (K (S I I)))
\end{lstlisting}

\textct{MABT} or whatever it is

Chris Barker
The iota combinator \textct{ι} can produce all the others, but in convoluted ways.

% https://stackoverflow.com/questions/7368507/expressing-y-in-term-of-ski-combinators-in-javascript

Combinator calculi are generally considered too minimalist for practical programming; they shine, however, in mathematical proofs.  This means that we can use them to define a standard of behavior for an algorithm, but implement the logic in a compliant form using whatever computational tools we like.  Thus it is better to think of Nock not as a bare programming language, but as a standard of computation which can under many circumstances be used as the computation itself.\footnote{Consideration of when to do so is largely driven by performance bottlenecks.}

\section{Nock Opcodes as Combinator}

The Nock specification is produced in full in Listing~\ref{lst:nock-4k}.

\begin{lstlisting}[style=listingcode,
                   caption={Nock 4K},
                   label={lst:nock-4k}]
Nock 4K

A noun is an atom or a cell. An atom is a natural number. A cell is an ordered pair of nouns.

Reduce by the first matching pattern; variables match any noun.

nock(a)             *a
[a b c]             [a [b c]]

?[a b]              0
?a                  1
+[a b]              +[a b]
+a                  1 + a
=[a a]              0
=[a b]              1

/[1 a]              a
/[2 a b]            a
/[3 a b]            b
/[(a + a) b]        /[2 /[a b]]
/[(a + a + 1) b]    /[3 /[a b]]
/a                  /a

#[1 a b]            a
#[(a + a) b c]      #[a [b /[(a + a + 1) c]] c]
#[(a + a + 1) b c]  #[a [/[(a + a) c] b] c]
#a                  #a

*[a [b c] d]        [*[a b c] *[a d]]

*[a 0 b]            /[b a]
*[a 1 b]            b
*[a 2 b c]          *[*[a b] *[a c]]
*[a 3 b]            ?*[a b]
*[a 4 b]            +*[a b]
*[a 5 b c]          =[*[a b] *[a c]]

*[a 6 b c d]        *[a *[[c d] 0 *[[2 3] 0 *[a 4 4 b]]]]
*[a 7 b c]          *[*[a b] c]
*[a 8 b c]          *[[*[a b] a] c]
*[a 9 b c]          *[*[a c] 2 [0 1] 0 b]
*[a 10 [b c] d]     #[b *[a c] *[a d]]

*[a 11 [b c] d]     *[[*[a c] *[a d]] 0 3]
*[a 11 b c]         *[a c]

*a                  *a
\end{lstlisting}

\subsection{\textct{I} Identity = Opcode 0}

Nock's binary tree addressing is fully general and yields the identity combinator \textct{I} as a special case.  (The cost, of course, is dragging in some machinery for basic arithmetic, as we will see.)

Conceptually, we could write this in a hybrid form between combinators and Nock thus:

\begin{lstlisting}[style=listingcode]
  *[a I b] == [0 1]
\end{lstlisting}

\subsection{\textct{K} Constant = Opcode 1}

\subsection{\textct{S} Substitution = Opcode 2}

\subsection{The Axiomatic Operators}

Because of Nock's cellular nature as a binary tree, we need the ability to access and manipulate nouns.

opcode 3, 4, 5, 10


\section{Idioms \& Design Patterns}

Several of the combinators we examined above are not present in Nock \emph{simpliciter}, but we can build them out of the other pieces.  There are also common idioms which are baked in as ``compound opcodes'' (such as conditional branching, opcode 6).

\subsection{\textct{C} Reversal}

C = S (S (K (S (K S) K)) S) (K K)

In Nock, the reversal combinator \textct{C} is almost free due to the binary tree nature of cells.

\begin{lstlisting}[style=listingcode]
  *[a C b] == *[7 [0 3] [[0 3] [0 2]]]
  *[a C b] == *[[0 7] [0 6]]
\end{lstlisting}


Recursion we get for free with opcode 0.

$\mu$-recursive functions


\section{Virtualization}

Nock is a crash-only pure language, which means at the base level it can only succeed and yield a new noun:  no side effects, no exceptions, no state.  Like a castaway on a desert island, we must build everything we need from raw materials at hand.

Nock has one trick built in and a couple of design patterns to help us get there:

\begin{enumerate}
  \item  Opcode 11 lets us send a hint to the runtime, which can be used to trigger side effects.  The Nock expression formally doesn't know anything about these, but we can use them to print results, read input, and so forth.
  \item  The runtime can maintain state outside of Nock and supply it as a subject to subsequent Nock expressions.  We will do a lot more with this later, but we want to highlight it as a key affordance now.
  \item  The virtualization pattern means that we can run Nock in Nock.
\end{enumerate}

Opcode 11 can encode anything we want:  the result of a hint is used by the runtime, and so anything the runtime is willing to work with can be produced.


The magic formula at the beating heart of a practical Nock operating system is simply:

\begin{lstlisting}[style=listingcode]
*([state -.events] [9 2 10 [6 0 3] 0 2])
\end{lstlisting}

\noindent
That is, apply the Nock formula against a subject pair consisting of the current state and the head of the list of events to process.

\renewcommand{\chaptertitle}{Interpreter Runtime}

%  Make first page footer:
\fancypagestyle{firststyle}{%
\fancyhf{}% Clear header/footer
\fancyhead{}
\fancyfoot[L]{{\footnotesize
              %% We toggle between these:
              % Manuscript submitted for review.\\
              \authorname{}~\authorpatp{}.  ``\chaptertitle{}''.  {\it The Nock Book} (2025):  10–$x$.
              }}
}
%  Arrange subsequent pages:
\fancyhf{}
\fancyhead[LE]{{\urbitfont The Nock Book}}
\fancyhead[RO]{\chaptertitle{}}
\fancyfoot[LE,RO]{\thepage}

\chapter{\chaptertitle}

\thispagestyle{firststyle}

\begin{abstract}

\end{abstract}

Now that we have established a firm foundation for idiomatic programming in Nock, we can build an interpreter using a higher-level language to execute Nock formulas.  At this point, we will build the simplest possible Nock interpreter, a tree-walking interpreter.  This requires a cursor and some basic memory management but not much more.  We will implement some minimalist static hints to permit side effects like output.

\renewcommand{\chaptertitle}{Moon:  A Higher-Level Language}

%  Make first page footer:
\fancypagestyle{firststyle}{%
\fancyhf{}% Clear header/footer
\fancyhead{}
\fancyfoot[L]{{\footnotesize
              %% We toggle between these:
              % Manuscript submitted for review.\\
              \authorname{}~\authorpatp{}.  ``\chaptertitle{}''.  {\it Nockonomicon} (2025):  40–$x$.
              }}
}
%  Arrange subsequent pages:
\fancyhf{}
\fancyhead[LE]{{\urbitfont Nockonomicon}}
\fancyhead[RO]{\chaptertitle{}}
\fancyfoot[LE,RO]{\thepage}

\chapter{\chaptertitle}

\thispagestyle{firststyle}

\begin{abstract}

\end{abstract}

While somewhat more economical than a minimal combinator language, Nock is still quite low-level and cumbersome as we have seen.  In this chapter, we will introduce a higher-level language called \emph{Moon} which compiles to a Nock noun.\footnote{The name \emph{Moon} is a double entendre:  it is a micro-\emph{Hoon}, the primary systems language of Urbit and NockApp applications, and it references Kleene's $\mu$-recursive functions.}

\section{Relationship of Moon to Nock}

You can simply understand our goal as to produce executable Nock nouns using a more legible and flexible syntax than pure Nock affords.  But it's worth casting Moon in the context of Nock's evaluation model.

The Moon compiler is a Nock noun which (as subject) operates on a text string to produce a new Nock noun.  This second noun is itself a formula against which well-structured formulas may be evaluated.  Formally, 

\begin{lstlisting}[style=listingcode]
Moon(source) == Nock(source moon-formula)
\end{lstlisting}

\noindent
that is, the evaluation of the Moon compiler as a subject against a source string produces a Nock noun which is the compiled output of the Moon compiler.  We can think of this as a two-step process:  first, we compile the source code to a Nock noun, and then we evaluate other Nock nouns against that noun.
% TODO check this statement closely, may be out of order


\noindent
where \textct{moon-source} is a text string containing the Moon source code, and \textct{moon-formula} is a Nock noun which is the compiled output of the Moon compiler.  The Moon compiler is itself a Nock noun, and so we can think of it as a self-compiling compiler.\footnote{We here follow the exposition of \citet{Yarvin2010c} in his description the a high-level language compiler to Nock for the first time.}

% Therefore, we need another deliverable: a human-usable language, Watt. Watt is defined as follows:
% Watt(source) == Nock(source watt-formula)
% Thus, for instance,
% urbit-formula == Watt(urbit-source) 
%               == Nock(urbit-source watt-formula)
% watt-formula == Watt(watt-source) 
%              == Nock(watt-source watt-formula)
% In English, Watt is a self-compiling compiler that translates a source file to a Nock formula. Watt is formally defined by the combination of watt-source and watt-formula. The latter verifies the former. Again, we are not in rocket-science territory here.

We will have to construct Moon in two stages:  a launch stage (``New Moon'') which can only compile simple expressions, a complex construction exemplifying all that we've seen; and a payload stage (``Crescent Moon'') which can compile itself and thus becomes a Nockhound's flywheel (``Full Moon'' being the completed product).  We have two possible ways to construct New Moon as the launch stage:  either we can write the Moon compiler directly in bespoke Nock by hand, or we can compose it in the same language as our interpreter, Rust.

\section{\textrm{🌑} New Moon}

\subsection{Design Criteria}

New Moon is complete when it can compile some set of simple expressions to Nock.  It does not yet, however, need to compile itself, which we defer until Crescent Moon.

Since Nock possesses its own set of idioms (as we've worked out in the preceding chapters), we know what the Nock products need to look like.  A higher-level language can take into account the structure of its target \textsc{isa} in its syntax or not, but we will find it easier to write Moon if we do.  (Hoon takes this road, being essentially a macro language and type system over Nock.)

For instance, in Nock we would apply a function by loading the reference via opcode 9 and then replacing its sample with the new arguments.  In Moon, we need a way to express this intent in an unambiguous way; let's say that we take a page out of Lisp's book and use parentheses to indicate function application.  Thus, we could write a hypothetical Moon expression like this:

\begin{lstlisting}[style=listingcode]
(+ 1 41)
\end{lstlisting}

\noindent
That's reasonable if Lisp-flavored; indeed, we will find that Moon has Lisp-like characteristics.  From a parsing standpoint, this needs to result in an abstract syntax tree (\textsc{ast}) which describes the operation and its arguments, thus detecting the head of the expression as its function and the tail to the ending parenthesis as arguments.  We'll use a syntax in which text is marked by a \textct{\%} cen prefix, resulting in the hypothetical \textsc{ast}:

\begin{lstlisting}[style=listingcode]
[%add 1 41]
\end{lstlisting}

\noindent
and ultimately in the Nock noun (modulo tree addresses):

\begin{lstlisting}[style=listingcode]
[8 [9 cor 0 arm] 9 2 10 [6 1 1 41] 0 2]
\end{lstlisting}

\noindent
(Read it carefully---we pin the core locally with an opcode 8, then we modify the local copy with opcode 10, replacing the sample with the new arguments.)

You can follow a thread of logic from the initial expression in Moon to the final Nock noun, mediated by the \textsc{ast} representation.  New Moon thus has two phases:  a parser and a compiler.  We can build them separately, but first let's look at more examples and their corresponding \textsc{ast} representations.

\subsection{New Moon Examples}

\begin{lstlisting}[style=listingcode]
::  This is a comment.  It doesn't affect the AST or code.

::  This is a simple expression.
(+ 1 41)

::  This is a more complex expression.
(- 54 (+ 3 9))

::  This is a variable declaration.
(= a 42)

::  This is a loop with a check.
for (= a 0)
=a (- )
\end{lstlisting}

\subsection{Parsing New Moon}

We will soon run out of symbols:  the punctuation in \textsc{ascii} is limited.  We could choose to switch to words throughout or to permit paired punctuation; we will do both eventually but for now we'll let it ride.\footnote{Another alternative is to support Unicode characters, as Julia does, but parsing \textsc{utf}-8 text is more complex than Moon will support out of the box.}


\subsection{Compiling New Moon}


\section{\textrm{🌙} Crescent Moon}

The objective of Crescent Moon is to bootstrap, or compile itself.  Once we have this capability, expanding Moon to a full language becomes simply a matter of adding more syntax and semantics to the compiler.  This means that Crescent Moon must be able to parse and compile text, with any required types being supplied as well.



At this point the question of the target \textsc{isa} becomes acute:  should Moon compile to strict Nock or to our extended Mock?  That is, if our compiler utilizes either extended opcodes 6--9 or Mock opcodes 12--XX, how will it relate to the runtime system?  At the end of the day, Nock is Nock (since we can always reduce), but we must make an aesthetic design choice which may have deep implications for Moon.

Let's consider what Hoon does.  Hoon actually targets Mock (with opcode 12 only) and does some sensible reductions and compile-time checks.  For instance, Hoon compilation of opcode 6 can eliminate branches if the conditional reduces to an opcode 1 (see \textct{+cond}).  The runtime compiler and system can make further optimizations:  opcode 11 \textct{\%fast} hints for jets are only checked for opcode 9 expressions, not for opcode 2.  Nock 7 is compiled away entirely.

At this point, we are not interested in optimizations, but there's no reason to forswear their future possibility.  We will therefore make the following design decision:  Moon targets canonical Nock without Mock opcodes (12 and higher).  This will make Moon somewhat less efficient than Hoon as a systems language, but it will remain more legible.



\section{\textrm{🌕} Full Moon}

From Crescent Moon, building Full Moon becomes a matter of expanding the language syntax, elaborating more parsers and other tools, and producing libraries of useful functions.

(at this point, back out from parentheses to centis for function application)

vase mode and type system redux

What other languages compile to Nock?  The earliest higher-level language was Watt, the ancestor of Hoon. %sss10.tex
Hoon
Jock



currying etc.
\input{txt/bootstrapping}
\input{txt/type_system}
\input{txt/event_loop}
\renewcommand{\chaptertitle}{A Virtual Machine Runtime}

%  Make first page footer:
\fancypagestyle{firststyle}{%
\fancyhf{}% Clear header/footer
\fancyhead{}
\fancyfoot[L]{{\footnotesize
              %% We toggle between these:
              % Manuscript submitted for review.\\
              \authorname{}~\authorpatp{}.  ``\chaptertitle{}''.  {\it Nockonomicon} (2025):  50–$x$.
              }}
}
%  Arrange subsequent pages:
\fancyhf{}
\fancyhead[LE]{{\urbitfont Nockonomicon}}
\fancyhead[RO]{\chaptertitle{}}
\fancyfoot[LE,RO]{\thepage}

\chapter{\chaptertitle}

\thispagestyle{firststyle}

\begin{quote}
\textbf{The Kraken}
Below the thunders of the upper deep,
Far, far beneath in the abysmal sea,
His ancient, dreamless, uninvaded sleep
The Kraken sleepeth: faintest sunlights flee
About his shadowy sides; above him swell
Huge sponges of millennial growth and height;
And far away into the sickly light,
From many a wondrous grot and secret cell
Unnumbered and enormous polypi
Winnow with giant arms the slumbering green.
There hath he lain for ages, and will lie
Battening upon huge sea worms in his sleep,
Until the latter fire shall heat the deep;
Then once by man and angels to be seen,
In roaring he shall rise and on the surface die.
(Alfred, Lord Tennyson)
\end{quote}

To date, we have run Nock either by hand or using a simple interpreter.  While this suffices for relatively small programs, our objective of building a Nock computer requires an ``operating function'', or persistent event loop.  This chapter introduces the considerations of running Nock continuously and reactively as well as how to maintain state and respond to opcode 11 hints.

\section{The Event Loop}

A single Nock expression as a subject--formula pair accepts two nouns and yields a new noun.  We can think of this new noun as the total state of a program or an operating system, feeding it into the next expression as the subject for a new formula to evaluate against.  This preserves the determinism and legibility of Nock while allowing for complex state transitions.

\subsection{A Simple State Machine}

A state machine is a mathematical abstraction for a system that has a definite state at any particular time, and which changes in response to inputs.

\subsubsection{Example:  Stoplight}

Think of a stoplight, which has a state of red, yellow, or green, and which changes its state in response to an external timer or a button press.


\subsubsection{Example:  Revolver}

Or again, of a six-shot revolver, which if we include chamber order has 64 finite state permutations with thirteen possible inputs---to ``load'' or ``clear'' at each chamber and to ``fire''---and two possible outputs---a ``bang'' or a ``click'' accompanied by a new state.

% 6 1
% 5 6
% 4 15
% 3 20
% 2 15
% 1 6
% 0 1

Aggregate or macro actions are also available---a ``moon clip'' corresponds to loading all six chambers at once, and a ``dump cylinder'' corresponds to clearing all six chambers at once.  The revolver is a state machine with a finite number of states, inputs, and outputs.

\footnote{\textct{+poke} and \textct{+peek} derive from \href{https://en.wikipedia.org/wiki/PEEK_and_POKE}{a very old convention in \textct{BASIC}}, wherein \textct{POKE} and \textct{PEEK} commands allowed programmers to write to and read from arbitrary memory addresses.}

We will call a state machine which adheres to the definition of its lifecycle function a \emph{kernel}.  Thus a kernel is a state machine of a certain shape which accepts nouns of a certain shape as input and produces nouns of a certain shape as output.  The kernel is the event loop and some wrapper of library functionality to bound and define the complete behavior of the state machine.  A kernel may be batch or continuous depending on whether its lifecycle function is reentrant.

\subsection{An Evolved State Machine}

While the two-armed model provides a straightforward way to manage state transitions, Nockhounds have found in practice that it is more convenient to supply four arms at the following axes:

\begin{enumerate}
  \item[\textct{+4}:] \textct{+load} is used in kernel upgrades, allowing the event loop to update itself live.
  \item[\textct{+10}:] \textct{+wish} accepts a core and parses it against a standard library for the kernel.
  \item[\textct{+22}:] \textct{+peek} grants read-only access to the system's state; this is often called a ``scry''.
  \item[\textct{+23}:] \textct{+poke} accepts events and processes them, resulting in a new state noun.
\end{enumerate}

\input{txt/userspace}
\input{txt/hints_and_jets}

\backmatter %%%%%%%%%%%%%%%%%%%%%%%%%%%%%%%%%%%%%%%%%%%%%%%%%%%%%%%%%%%%%%%%%%%%

\appendix
\input{txt/hoon_proper}

% 1. The roots of computing 
% 2. Practical combinators & design patterns/idioms  (including virtualization)
% 3. An interpreter runtime 
% 4. A higher level language (micro Hoon, Moon)
% 5. Bootstrapping Moon (parsers etc) 
% 6. Types, doors, etc.
% 7. The event loop
% 8. A toy runtime
% 9. Userspace loops upon loops, scrying
% 10. Jets and hints
% 11. Hoon proper, Jock

\selectlanguage{USenglish}
\printbibliography
\end{document}
